\chapter{Environmental Pollution, Control Technologies and Waste Management}
\label{chap:pollution-waste}

\section{Overview and Course Learning Outcomes}
Environmental pollution represents the contamination of physical environmental media (air, water, soil) through the introduction of chemical xenobiotics, particulates, noise, or ionizing radiation at levels exceeding the natural assimilative capacity of recipient ecosystems. This chapter examines ambient air quality standards, point-source industrial emission controls, aquatic oxygen sag kinetics (Streeter-Phelps), soil pollution, noise abatement, emerging micro-pollutants, and modern solid and biomedical waste management frameworks.

\begin{important}[Course Learning Outcomes for Unit 4]
Upon completing this chapter, the student will be capable of:
\begin{enumerate}[leftmargin=1.5em]
    \item Evaluating criteria air pollutants against CPCB National Ambient Air Quality Standards (NAAQS) and calculating Air Quality Indices (AQI).
    \item Analyzing air pollution control equipment: settling chambers, cyclones, venturi scrubbers, and electrostatic precipitators (ESPs).
    \item Deriving and solving the Streeter-Phelps dissolved oxygen deficit equation for point-source river pollution.
    \item Evaluating ambient noise levels in dB(A) against the Noise Pollution (Regulation and Control) Rules, 2000.
    \item Designing integrated municipal solid waste management and biomedical waste segregation schemes under Indian statutory rules.
\end{enumerate}
\end{important}

\section{Air Pollution and Particulate Control Technologies}
Air pollutants are classified into primary pollutants emitted directly from source points ($CO$, $SO_2$, $NO_x$, primary $\text{PM}_{10}/\text{PM}_{2.5}$) and secondary pollutants formed through atmospheric photochemical reactions ($O_3$, peroxyacetyl nitrate --- PAN, secondary sulfates and nitrates).

\subsection{Particulate Removal Equipment Mechanics}
\begin{itemize}[leftmargin=1.5em]
    \item \textbf{Gravitational Settling Chambers:} Relies on gravity settling ($v_t = \frac{g d_p^2 (\rho_p - \rho)}{18 \mu}$). Effective only for large particulates ($d_p > 50\ \mu\text{m}$).
    \item \textbf{Cyclone Separators:} Employs centrifugal force generated by vortex airflow. High separation efficiency for coarse particles ($d_p > 10\ \mu\text{m}$).
    \item \textbf{Electrostatic Precipitators (ESPs):} Flue gas traverses high-voltage corona discharge electrodes ($-40\text{ to }-70\text{ kV}$), imparting negative charge to particulates, which then migrate electrostatically to grounded collecting plates. Collection efficiency exceeding $99.5\%$ for fine fly ash ($d_p \ge 0.1\ \mu\text{m}$).
    \item \textbf{Fabric Filter Baghouses:} Flue gas passes through woven fabric filter bags; achieves $>99.9\%$ efficiency across all particle sizes via inertial impaction, interception, and Brownian diffusion.
    \item \textbf{Wet Scrubbers (Venturi):} Employs high-velocity liquid atomization in a converging-diverging venturi throat to capture fine particulates and absorb soluble acid gases ($SO_2, HCl$).
\end{itemize}

\section{Water Pollution and Oxygen Sag Modeling}
Point-source discharge of biodegradable organic waste exerts a biochemical oxygen demand on stream dissolved oxygen (DO).

\begin{keyequation}[Streeter-Phelps Dissolved Oxygen Deficit Equation]
The spatial and temporal variation of DO deficit $D(t) = \text{DO}_{\text{sat}} - \text{DO}(t)$ in a flowing stream receiving wastewater is governed by simultaneous deoxygenation and reaeration:
\begin{equation}
\frac{dD}{dt} = k_d L - k_r D
\end{equation}
Integrating under initial boundary conditions $L(0) = L_0$ and $D(0) = D_0$:
\begin{equation}
D(t) = \frac{k_d L_0}{k_r - k_d} \left( e^{-k_d t} - e^{-k_r t} \right) + D_0 e^{-k_r t}
\end{equation}
where $k_d$ is the deoxygenation rate constant ($\text{day}^{-1}$), $k_r$ is the reaeration rate constant ($\text{day}^{-1}$), $L_0$ is the initial ultimate BOD ($\text{mg/L}$), and $D_0$ is the initial oxygen deficit ($\text{mg/L}$).
\end{keyequation}

\begin{example}[Critical Oxygen Deficit Calculation]
A treated wastewater discharge enters a river producing initial mixed $L_0 = 25.0\text{ mg/L}$ and initial deficit $D_0 = 1.20\text{ mg/L}$. The laboratory deoxygenation constant is $k_d = 0.20\text{ day}^{-1}$, stream reaeration constant is $k_r = 0.50\text{ day}^{-1}$, and saturation DO is $\text{DO}_{\text{sat}} = 9.0\text{ mg/L}$. Determine the critical time $t_c$, the critical deficit $D_c$, and the minimum DO concentration.
\end{example}

\begin{solutionbox}
Using the critical time relation ($k_r - k_d = 0.30$):
\begin{align*}
t_c &= \frac{1}{k_r - k_d} \ln\left\{ \frac{k_r}{k_d} \left[ 1 - D_0 \frac{k_r - k_d}{k_d L_0} \right] \right\} \\
&= \frac{1}{0.30} \ln\left\{ \frac{0.50}{0.20} \left[ 1 - 1.20 \times \frac{0.30}{0.20 \times 25.0} \right] \right\} \\
&= \frac{1}{0.30} \ln\left\{ 2.50 \times \left[ 1 - \frac{0.36}{5.0} \right] \right\} = \frac{1}{0.30} \ln(2.50 \times 0.928) = \frac{1}{0.30} \ln(2.32) \\
&= \frac{0.84157}{0.30} = 2.805\text{ days}.
\end{align*}
At $t = t_c$, deoxygenation rate equals reaeration rate ($k_r D_c = k_d L(t_c)$):
\begin{align*}
D_c &= \frac{k_d}{k_r} L_0 e^{-k_d t_c} = \frac{0.20}{0.50} \times 25.0 \times e^{-0.20 \times 2.805} = 10.0 \times e^{-0.561} \\
&= 10.0 \times 0.5706 = 5.71\text{ mg/L}.
\end{align*}
The minimum dissolved oxygen in the river is:
\begin{equation*}
\text{DO}_{\text{min}} = \text{DO}_{\text{sat}} - D_c = 9.0 - 5.71 = 3.29\text{ mg/L}.
\end{equation*}
Because $3.29\text{ mg/L} < 4.0\text{ mg/L}$ (CPCB minimum standard for fish survival), advanced biological treatment is required before effluent discharge.
\end{solutionbox}

\section{Solid and Biomedical Waste Management}
Modern waste management prioritizes the circular economy waste hierarchy:
\textbf{Prevention $\to$ Minimization $\to$ Reuse $\to$ Recycling $\to$ Energy Recovery (Waste-to-Energy) $\to$ Sanitary Landfill Disposal}.

\subsection{Biomedical Waste Management Rules, 2016}
Segregation at source is mandatory using standardized color-coded receptacles:
\begin{itemize}[leftmargin=1.5em]
    \item \textbf{Yellow Container (Non-Chlorinated Plastic):} Human anatomical tissue, animal carcasses, soiled cotton, expired pharmaceuticals. Disposal: High-temperature incineration ($>1100^\circ\text{C}$ secondary chamber) or plasma pyrolysis.
    \item \textbf{Red Container:} Contaminated recyclable plastics (catheters, IV bottles, tubing, syringes without needles). Disposal: Autoclaving/Microwaving followed by mechanical shredding and plastic recycling.
    \item \textbf{White Translucent (Puncture-Proof):} Sharps, disposable needles, scalpels, surgical blades. Disposal: Autoclaving followed by encapsulation in a concrete sharp pit.
    \item \textbf{Blue Box / Cardboard:} Glass ampoules, medicine vials, metallic orthopedic implants. Disposal: Chemical disinfection with sodium hypochlorite followed by autoclaving.
\end{itemize}

\begin{quickrevision}[Chapter 4 Summary]
\begin{itemize}[leftmargin=1.2em]
    \item Electrostatic Precipitators (ESPs) achieve $>99.5\%$ efficiency for fine fly ash particulates.
    \item Streeter-Phelps oxygen sag models deoxygenation ($k_d$) and atmospheric reaeration ($k_r$).
    \item CPCB requires minimum $\text{DO} \ge 4.0\text{ mg/L}$ to sustain aquatic vertebrate life.
    \item Noise rules prescribe residential limits of $55\text{ dB(A)}$ (Day) and $45\text{ dB(A)}$ (Night).
    \item Biomedical waste is strictly segregated into Yellow, Red, White, and Blue color streams.
\end{itemize}
\end{quickrevision}
